\section{Publish the state and frame at one boundary}
A simulation step updates both axes, then renders the resulting position. The
walkthrough also renders the initial state before taking any steps. For a request
of \(N\) steps, the observable sequence therefore has \(N+1\) frames, numbered
zero through \(N\). A request of zero still publishes the initial frame.

At each stopped \texttt{frame\_ready} boundary, publish the four raw state words,
the tick number, the selected pixel ID, and the complete frame. A debugger may
stop there before the supplied marker instruction executes. This gives the
numeric and visual observations one meaning: every word belongs to the same
simulation state.

The initial inputs, requested step count, and rounding mode are read-only. Guards
and a spare caller-owned frame remain untouched by the main animation. The main
frame starts at the configured Bitmap Display base. An independently called
renderer instead receives its destination pointer as an argument, so it can
render into the spare frame or a relocated allocation as well.

The frame represents the current state, not a trail. Clearing the previous white
cell is part of producing the next complete frame. The reference can clear all
64 words and then draw the new point; an alternative can classify every cell in
one pass. Correctness at the boundary does not prescribe either implementation
or an exact instruction count. Preserve the calling convention when the main
loop calls the axis updater and renderer: full saved-register contents and the
stack pointer survive each call.

The supplied single-instruction marker between \texttt{frame\_ready} and
\texttt{frame\_continue} also provides a countable publication boundary. Observing
that region distinguishes the expected initial-and-step sequence from an extra
or omitted frame. The marker remains part of the teaching harness. It is not a
delay instruction for real time, and it does not make a simulator step last a
specified number of milliseconds.
